\documentclass[11pt]{amsart}

\usepackage[margin=1in]{geometry}
\usepackage{amsmath,amssymb,mathtools}
\usepackage[T1]{fontenc}
\usepackage{lmodern}
\usepackage{microtype}
\usepackage{xcolor}
\usepackage{xurl}
\usepackage[colorlinks=true,
            linkcolor=blue!55!black,
            citecolor=blue!55!black,
            urlcolor=blue!55!black]{hyperref}

\newtheorem{theorem}{Theorem}[section]
\newtheorem{proposition}[theorem]{Proposition}
\newtheorem{lemma}[theorem]{Lemma}
\newtheorem{corollary}[theorem]{Corollary}
\theoremstyle{remark}
\newtheorem{remark}[theorem]{Remark}

\newcommand{\Z}{\mathbf Z}
\newcommand{\F}{\mathbf F}
\newcommand{\Ga}{\mathbb G_a}
\newcommand{\Gm}{\mathbb G_m}
\newcommand{\Spec}{\operatorname{Spec}}
\newcommand{\Soc}{\operatorname{Soc}}
\newcommand{\eps}{\varepsilon}
\newcommand{\m}{\mathfrak m}
\newcommand{\ot}{\otimes}
\newcommand{\unit}{\mathrm e}

\title[A rank-four counterexample]
      {A Rank-Four Counterexample to\\
       Grothendieck's Power Question}
\author{}
\date{}
\subjclass[2020]{14L15, 16T05}
\keywords{finite flat group scheme, power word, Oort--Tate group scheme,
matched product, affine group}

\begin{document}

\begin{abstract}
The power question customarily attributed to Grothendieck asks whether a
finite locally free group scheme of rank \(n\) is killed by its
\(n\)-th power word.  We give a counterexample of rank four.  Let
\[
 R=\Z[a,b]/(a^3,b^3,a^2b+2)
\]
and
\[
 A=R[U,V]/(U^2-abU+b^2V,\ V^2-a^2V).
\]
Then \(A\) is free of rank four over \(R\).  If
\(\lambda=(1+aU)(1+bV)\), the formulas
\[
 \Delta(U)=U\ot1+\lambda\ot U,\qquad
 \Delta(V)=V\ot\lambda+1\ot V
\]
make \(A\) a Hopf algebra, but
\[
 [4]^\#(U)=2bUV\ne0,\qquad [4]^\#(V)=0,\qquad [8]=\unit.
\]

The construction takes place inside
\(H=(\Ga^2)\rtimes\Gm\).  Modulo \(q=b^2\), it is an exact matched
product of two order-two Oort--Tate groups.  A general construction of
such rank-four matched products is given.  Over \(R\), one scalar
compatibility fails by \(aq\); replacing
\(U^2-abU=0\) by \(U^2-abU=-qV\) restores closure under multiplication.
The resulting group is nevertheless not killed by four: the relevant
coefficient is \(a^2q=-2b\ne0\).
The group structure follows from the group structure on the closed
fiber.  A companion Macaulay2 file checks all polynomial identities.
\end{abstract}

\maketitle

\begin{center}
\fbox{\begin{minipage}{0.92\textwidth}
\small
\textbf{Disclosure: AI-generated work.}\ This document in its entirety---the
construction, the proofs, the computations, and the exposition---was generated
by artificial intelligence. The work was carried out by the AI models Codex
(OpenAI) and Claude (Anthropic).
\end{minipage}}
\end{center}
\medskip

\section{Introduction and history}

Let \(G/S\) be a finite locally free group scheme of constant rank \(n\).
The multiplication and diagonal define the power word
\[
 [n]\colon G\longrightarrow G,\qquad g\longmapsto g^n.
\]
If \(G\) is not commutative, this need not be a homomorphism, but it is
still a morphism of \(S\)-schemes.  The natural analogue of Lagrange's
theorem asks whether it is the constant unit word
\[
 \unit_G:G\longrightarrow S\stackrel{\unit}{\longrightarrow}G.
\]

The question is traditionally attributed to Grothendieck and has a long,
unusually explicit history.  SGA~3 proves that a finite group scheme over
a field is killed by its rank
\cite[Exp.~VIIA, Prop.~8.5]{SGA3I}.  The field result immediately
implies the result over a reduced base: the difference between the
comorphisms of \([n]\) and the constant unit word vanishes after passage
to every residue field.  In a local free basis its coefficients
therefore lie in every prime ideal, and hence vanish when the base is
reduced.

There is a bibliographic wrinkle in the usual attribution.  In the
printed 1970 version of Expos\'e~VIII, Lemma~7.3 was stated for a
commutative group scheme, so the adjacent Remark~7.3.1 literally asked
for the arbitrary-base result in that context
\cite[Exp.~VIII, Lem.~7.3 and Rem.~7.3.1]{SGA3II}.  The modern
annotated edition removes the commutativity hypothesis and explicitly says that it
has modified the remark accordingly
\cite[notes (38)--(39)]{SGA3VIIIannotated}.  Tate and Oort nevertheless
formulated the arbitrary-base \emph{noncommutative} question explicitly
in 1970 \cite[p.~5]{OortTate}, and Schoof and Torti subsequently
formulated it as Grothendieck's question.  We follow that established
terminology.  Tate's 1997 survey also records the problem as open
\cite[\S3.8]{Tate1997}.

Deligne found a proof for every \emph{commutative} finite locally free
group scheme over an arbitrary base.  He communicated the proof to Tate
and Oort, who published it in 1970 in their paper on group schemes of
prime order \cite[\S1]{OortTate}.  Thus any counterexample must use both a
nonreduced base and a noncommutative group law.

Oort and Tate also proved that every group scheme of prime rank \(p\)
over an arbitrary base is commutative and killed by \(p\)
\cite[Thm.~1]{OortTate}.  Hence ranks two and three are impossible, and
rank four is the first rank at which a counterexample could occur.

Schoof established a principal general nilpotence criterion over
nonreduced bases.  If \(R\) is a local Artin ring with residue
characteristic \(p\), maximal ideal \(\m_R\), and
\[
 \m_R^p=p\m_R=0,
\]
then every finite flat group scheme over \(R\) is killed by its rank
\cite{Schoof2001}.  Schoof later gave a concise account of the problem,
the reductions to local Artin bases, the field case, and Deligne's
argument \cite{SchoofSurvey}.  Schoof also proved, without a length
restriction, killedness for deformations of the standard noncommutative
matrix group \(\alpha_p\rtimes\mu_{p^m}\).  More recently, Torti
extended this deformation-theoretic analysis to an explicit family
\(\alpha_p\rtimes_\lambda\mu_{p^m}\) of noncommutative special fibers and
proved killedness for all of them without a bound on the length of the
base \cite[Thms.~1.3--1.4]{Torti}.  That work, published in December
2025, still presents the general question as open.  We know of no earlier
counterexample.

The example below gives a negative answer.  It lies just outside the
known positive ranges.  Its residue characteristic is two, but
\[
 \m^2\ne0,\qquad 2\m\ne0.
\]
Its special fiber is the \emph{commutative} group
\(\alpha_2\times\alpha_2\); noncommutativity itself appears only in the
nilpotent deformation.  In particular, the example is not a deformation
of one of the noncommutative special fibers treated by Schoof and Torti.

The preceding positive results also suggest what a counterexample would
have to look like.  Nilpotents in the base are essential, since the
assertion holds over every reduced base.  Noncommutativity is essential,
since Deligne's theorem covers the commutative case.  Yet it is not
enough simply to deform a familiar noncommutative group in the most
elementary way, as the results of Schoof and Torti show.  Our example
uses a different possibility: the special fiber is commutative, and a
noncommutative interaction is created entirely by nilpotent
coefficients.  The base retains enough of that interaction to repair
multiplicative closure while leaving a nonzero trace in the power word.

Here is the complete statement.

\begin{theorem}\label{thm:main}
Let
\begin{equation}\label{eq:R}
 R=\Z[a,b]/(a^3,b^3,a^2b+2)
\end{equation}
and
\begin{equation}\label{eq:A}
 A=R[U,V]/(U^2-abU+b^2V,\ V^2-a^2V).
\end{equation}
Put \(W=UV\) and
\begin{equation}\label{eq:lambda}
 \lambda=(1+aU)(1+bV).
\end{equation}
Then:
\begin{enumerate}
\item \(R\) is a local Artin complete intersection of length nine,
with residue field \(\F_2\) and characteristic four.
\item \(A\) is free of rank four over \(R\), on \(1,U,V,W\).
\item There is a Hopf algebra structure on \(A\) with
\[
 \eps(U)=\eps(V)=0,\qquad
 \Delta(U)=U\ot1+\lambda\ot U,\qquad
 \Delta(V)=V\ot\lambda+1\ot V.
\]
\item For \(G=\Spec A\),
\[
 [4]^\#(U)=2bW\ne0,\qquad
 [4]^\#(V)=0,\qquad
 [8]=\unit.
\]
\item The special fiber is \(\alpha_2\times\alpha_2\), whereas \(G\)
itself is noncommutative.
\end{enumerate}
In particular, \(G\) is a finite locally free group scheme of rank four
which is not killed by four.
\end{theorem}

The proof is organized by the ambient affine group
\(H=(\Ga^2)\rtimes\Gm\).  Its triangular-matrix representation
makes associativity and the power words formal.  Section~3.1 constructs
rank-four subgroups of \(H\) as matched products of two order-two
Oort--Tate groups and identifies the quotient \(G_0\).  Section~3.2
shows that one cross compatibility fails over \(R\), and that the
single correction \(qV\) restores multiplicative closure.  The power
calculation is then reduced to a geometric sum in Section~4.
\section{The base and the affine-group mechanism}

We begin with the small amount of commutative algebra needed to detect
the final nonzero class.

\begin{lemma}\label{lem:base}
The additive group of \(R\) has the decomposition
\begin{equation}\label{eq:additive-R}
 R_{\mathrm{add}}\simeq
 (\Z/4)\{1,b\}\oplus
 (\Z/2)\{a,a^2,ab,b^2,ab^2\}.
\end{equation}
In particular,
\[
 4=0,\qquad 2b=-a^2b^2\ne0.
\]
Moreover \(R\) is a local Artin complete intersection of length nine,
with maximal ideal \(\m=(a,b)\) and residue field \(\F_2\).
\end{lemma}

\begin{proof}
Let \(R_0=\Z[a,b]/(a^3,b^3)\), free over \(\Z\) on the nine box
monomials \(a^ib^j\), \(0\le i,j\le2\).  Multiplication by \(a^2b\)
has only two nonzero arrows:
\[
 1\longmapsto a^2b,\qquad b\longmapsto a^2b^2.
\]
The additive group of \(R\) is the cokernel of
\(2+\text{multiplication by }a^2b\).  On the ordered pairs
\[
 (1,a^2b),\qquad (b,a^2b^2),
\]
this endomorphism has matrix
\[
 \begin{pmatrix}2&0\\1&2\end{pmatrix},
\]
whose Smith normal form is \(\operatorname{diag}(1,4)\).  On each of
the remaining monomials
\[
 a,\ a^2,\ ab,\ b^2,\ ab^2
\]
it is multiplication by \(2\).  The two \(2\times2\) blocks therefore
give cyclic summands generated by \(1\) and \(b\), each of order four,
while the five \(1\times1\) blocks give the five displayed summands of
order two.  This proves \eqref{eq:additive-R}.  It also shows directly
that \(2b=-a^2b^2\) is nonzero.

Since \(a\) and \(b\) are nilpotent, they belong to every prime ideal
of \(R\).  The relation \(2=-a^2b\) then implies that every prime also
contains \(2\).  Thus \((a,b)\) is the unique prime, hence the unique
maximal ideal, and \(R/(a,b)=\F_2\).  The displayed additive
decomposition has order \(2^9\), so the length over the residue field
is nine.  It shows that \(4=0\) but \(2\ne0\), and therefore that
\(\operatorname{char}(R)=4\).

Finally, \(a^3,b^3\) form a regular sequence in \(\Z[a,b]\).
Multiplication by \(a^2b+2\) on the free abelian group \(R_0\) has
determinant
\[
 4^2\cdot2^5=2^9
\]
by the same block decomposition.  It is therefore injective.  Hence
\(a^3,b^3,a^2b+2\) is a regular sequence, proving the
complete-intersection assertion.
\end{proof}

We will repeatedly use the immediate consequences
\begin{equation}\label{eq:base-identities}
 2a=4=0,\qquad
 q:=b^2,\qquad
 q^2=bq=2q=0,\qquad
 a^2q=-2b.
\end{equation}
The first equality follows by multiplying \(a^2b=-2\) by \(a\), and the
others follow from \(b^3=0\).

The ideal \(Q=(q)\) is square-zero.  The following sequence of
nonvanishing products will be used twice:
\begin{equation}\label{eq:coefficient-chain}
 q\longmapsto aq\longmapsto a^2q=-2b\ne0,
 \qquad a^3q=0.
\end{equation}
Here \(q=b^2\), \(aq=ab^2\), and \(a^2q=-2b\) are all nonzero by
\eqref{eq:additive-R}.
The first product \(aq\) is the obstruction to lifting the matched
product of Section~3.1 without changing its equations.  The correction
by \(q\) removes that obstruction, while the second product \(a^2q\)
reappears in the fourth-power word.

Now consider the affine group
\begin{equation}\label{eq:H}
 H=\Spec R[u,v,z,z^{-1}]
\end{equation}
with multiplication
\begin{equation}\label{eq:H-law}
 (u,v,z)(u',v',z')
   =(u+zu',\,vz'+v',\,zz').
\end{equation}
Thus, for every \(R\)-algebra \(C\), an element of \(H(C)\) is a triple
\((u,v,z)\in C\times C\times C^\times\).  The coordinate \(z\) is a
multiplicative scale.  It acts on the next upper-translation coordinate
from the left and on the preceding lower-translation coordinate from
the right.  The inverse is
\[
 (u,v,z)^{-1}=(-z^{-1}u,-vz^{-1},z^{-1}).
\]
It is the fiber product over \(\Gm\) of an upper and a lower affine
group, as shown by the faithful representation
\begin{equation}\label{eq:H-matrix}
 (u,v,z)\longmapsto
 \begin{pmatrix}z&u\\0&1\end{pmatrix}
 \oplus
 \begin{pmatrix}z&0\\v&1\end{pmatrix}.
\end{equation}
After the change of coordinate \(v=zy\), this is the usual semidirect
product \(\Ga^2\rtimes\Gm\), with weights \(1\) and \(-1\).

On coordinate rings, the same law says
\[
 \Delta(u)=u\ot1+z\ot u,\qquad
 \Delta(v)=v\ot z+1\ot v,\qquad
 \Delta(z)=z\ot z.
\]
The last identity says geometrically that \(z:H\to\Gm\) is a group
homomorphism; this is what it means for \(z\) to be \emph{group-like}.
The first two identities say that \(u\) and \(v\) are translation
cocycles twisted by that character; algebraically they are respectively
left and right \(z\)-skew primitive.  When \(z=1\), they become ordinary
primitive coordinates, hence homomorphisms to \(\Ga\).  This dictionary
is the reason for using \(H\): once a closed subscheme has \(z=\lambda\)
and is stable under the displayed multiplication, the desired Hopf
formulas and all their associativity identities come for free.

The matrix description contains the power calculation we need.  Write
\[
 N_n(T)=1+T+\cdots+T^{n-1}.
\]
Then
\begin{equation}\label{eq:H-power}
 (u,v,z)^n=\bigl(N_n(z)u,\ vN_n(z),\ z^n\bigr).
\end{equation}
This is just the geometric-series formula for the off-diagonal entry of
a triangular matrix power.  It replaces an iterated Hopf-algebra
calculation by a single scalar \(N_n(z)\).
Equivalently, a function \(U\) satisfying
\[
 \Delta(U)=U\ot1+\lambda\ot U
\]
is a translation coordinate in an affine-group representation, and
\([n]^\#(U)=N_n(\lambda)U\).  This is the geometric meaning of a
\(\lambda\)-skew primitive element.

We shall use the following standard lifting observation to pass from
monoids to groups.

\begin{lemma}\label{lem:monoid-to-group}
Let \(S\) be a local Artin ring and let \(M\) be a finite locally free
monoid scheme over \(S\).  If the closed fiber of \(M\) is a group
scheme, then \(M\) is a group scheme.
\end{lemma}

\begin{proof}
Let \(B=\mathcal O(M)\).  In the convolution algebra
\(\operatorname{End}_S(B)\), the reduction of
\(\mathrm{id}_B\) has an inverse, namely the antipode of the closed
fiber.  Since \(B\) is finite free, lift that antipode to an
\(S\)-linear endomorphism \(s\) of \(B\).  Write \(*\) for convolution
and \(e=\eta\eps\) for its unit.  Both
\[
 a=s*\mathrm{id}_B,\qquad b=\mathrm{id}_B*s
\]
belong to \(e+\mathfrak m_S\operatorname{End}_S(B)\).  The latter is a
nilpotent ideal in the convolution algebra, so the finite geometric
series makes \(a\) and \(b\) invertible.  Now
\[
 s_{\mathrm L}=a^{-1}*s,\qquad
 s_{\mathrm R}=s*b^{-1}
\]
satisfy
\[
 s_{\mathrm L}*\mathrm{id}_B=e,\qquad
 \mathrm{id}_B*s_{\mathrm R}=e.
\]
Consequently
\[
 s_{\mathrm L}
 =s_{\mathrm L}*(\mathrm{id}_B*s_{\mathrm R})
 =(s_{\mathrm L}*\mathrm{id}_B)*s_{\mathrm R}
 =s_{\mathrm R}.
\]
This common endomorphism is an antipode for \(B\).
\end{proof}

\section{Oort--Tate matched products and the corrected lift}
\label{sec:construction}

Put
\[
 q=b^2,\qquad Q=(q)\subset R,\qquad D=R/Q.
\]
Thus \(Q^2=0\).  We first construct the group modulo \(Q\) from two
order-two Oort--Tate groups.

\subsection{An Oort--Tate matched product downstairs}

Let \(S\) be a ring, and let \(c,d\in S\) satisfy \(cd=-2\).  In the
sign convention used here, the order-two Oort--Tate group is
\begin{equation}\label{eq:OT2}
 \operatorname{OT}_S(c,d)
 =
 \Spec S[T]/(T^2-cT),
 \qquad
 \Delta(T)
 =T\ot1+1\ot T+dT\ot T.
\end{equation}
Equivalently, if \(\chi=1+dT\), then
\[
 \Delta(T)=T\ot1+\chi\ot T,
 \qquad
 \Delta(\chi)=\chi\ot\chi,
 \qquad
 \chi^2=1.
\]
These identities also give a self-contained verification of the Hopf
structure.  Indeed,
\[
 \chi^2
 =1+(2d+cd^2)T
 =1+d(2+cd)T
 =1,
\]
and the formula for \(\Delta(\chi)\) follows by expanding
\(1+d\Delta(T)\).  If
\[
 T_*=T_1+(1+dT_1)T_2,
\]
then, using \(T_i^2=cT_i\) and \((1+dT_1)^2=1\), one finds
\[
 T_*^2-cT_*
 =
 \bigl(c(1-(1+dT_1))+2(1+dT_1)T_1\bigr)T_2
 =0.
\]
The last equality is exactly \(cd=-2\).
Thus \(\Delta(T)\) satisfies the same quadratic relation as \(T\), so
\(\Delta\) descends to the quotient.  Coassociativity follows without
another calculation: both iterated coproducts of \(T\) equal
\[
 T\ot1\ot1+\chi\ot T\ot1+\chi\ot\chi\ot T.
\]
The counit is \(\eps(T)=0\), and the antipode is
\(S(T)=T\), since
\[
 T+\chi T=2T+dT^2=(2+dc)T=0.
\]
Thus \(\operatorname{OT}_S(c,d)\) is finite locally free of rank two
and every point has order dividing two.
This is the \(p=2\) case of the Oort--Tate classification
\cite[Introduction]{OortTate}; our parameter \(d\) is the negative of
the second parameter there, so their relation \(c(-d)=2\) becomes
\(cd=-2\).  Over a local base, every order-two group scheme with free
augmentation ideal has this form.

The following criterion gives a rank-four analogue adapted to the
ambient group \(H\).

\begin{proposition}[Matched products of two Oort--Tate groups]
\label{prop:OT-matched}
Let \(c,d,e,f\in S\) satisfy \(cd=ef=-2\), and set
\[
 A_{c,d;e,f}
 =
 S[U,V]/(U^2-cU,\ V^2-eV),
 \qquad
 \lambda=(1+dU)(1+fV).
\]
Then
\begin{equation}\label{eq:OT-matched-graph}
 \Spec A_{c,d;e,f}\longrightarrow H_S,
 \qquad
 (U,V)\longmapsto(U,V,\lambda)
\end{equation}
is a subgroup of \(H_S\) if and only if
\begin{equation}\label{eq:OT-matched-conditions}
 cf=ed=0.
\end{equation}
When these conditions hold, it is finite locally free of rank four and
its underlying
scheme is
\[
 \operatorname{OT}_S(c,d)
 \times_S
 \operatorname{OT}_S(e,f).
\]
If \(p(u)=(u,0,1+du)\) and \(k(v)=(0,v,1+fv)\), the two factors satisfy
\begin{equation}\label{eq:OT-refactor}
 k(v)p(u)
 =
 p\bigl((1+fv)u\bigr)
 k\bigl(v(1+du)\bigr).
\end{equation}
\end{proposition}

\begin{proof}
Assume first that \(cf=ed=0\).
Regard the quotient first as
\[
 \bigl(S[V]/(V^2-eV)\bigr)[U]/(U^2-cU).
\]
The two successive monic relations give the free basis \(1,U,V,UV\).
Put
\(L=1+dU\) and \(M=1+fV\).  Since \(cd=ef=-2\),
\[
 L^2=M^2=1;
\]
hence \(\lambda=LM\) is a unit.  The comorphism
\[
 S[u,v,z,z^{-1}]\longrightarrow A_{c,d;e,f},
 \qquad
 (u,v,z)\longmapsto(U,V,\lambda)
\]
sends \(z^{-1}\) to \(\lambda^{-1}\) and is surjective because its
image contains \(U\) and \(V\).  Thus \eqref{eq:OT-matched-graph} is a
finite locally free closed immersion of rank four.

It remains to check that the product of the two Oort--Tate factors is
stable under multiplication.  In \(H_S\),
\[
 k(v)p(u)
 =
 \bigl((1+fv)u,\,v(1+du),\,(1+fv)(1+du)\bigr).
\]
Write
\[
 \widetilde u=(1+fv)u,
 \qquad
 \widetilde v=v(1+du).
\]
The cross conditions in \eqref{eq:OT-matched-conditions} give
\[
 \widetilde u^{\,2}-c\widetilde u=-cfuv=0,
 \qquad
 \widetilde v^{\,2}-e\widetilde v=-eduv=0.
\]
Thus \(\widetilde u\) and \(\widetilde v\) lie in the two factors.

We also need their characters to multiply to the ambient third
coordinate.  Put \(\alpha=du\), \(\beta=fv\), and
\(E=\alpha\beta\).  A direct expansion gives
\[
 (1+d\widetilde u)(1+f\widetilde v)
 -(1+\beta)(1+\alpha)
 =
 E(2+\alpha+\beta+E).
\]
The scalar relations imply \(2d=2f=0\): multiply \(ed=0\) by \(f\),
and \(cf=0\) by \(d\).  Using \(u^2=cu\) and
\(v^2=ev\), one then obtains
\[
 2E=E\alpha=E\beta=E^2=0.
\]
For clarity, the middle two equalities are
\[
 E\alpha=d^2f u^2v
          =cdfd\,uv=-2df\,uv=0,
 \qquad
 E\beta=df^2uv^2
          =def^2\,uv=-2df\,uv=0,
\]
and \(E^2=(E\alpha)\beta=0\).
The character identity follows, proving the refactorization
\eqref{eq:OT-refactor}.

Every point of the graph has a unique factorization \(p(u)k(v)\).
To multiply two such points, use \eqref{eq:OT-refactor} on the middle
two factors:
\[
 p(u_1)k(v_1)p(u_2)k(v_2)
 =
 p(u_1)\,
 p\bigl((1+fv_1)u_2\bigr)\,
 k\bigl(v_1(1+du_2)\bigr)\,
 k(v_2).
\]
The two \(p\)-terms multiply in \(\operatorname{OT}_S(c,d)\), and the
two \(k\)-terms multiply in \(\operatorname{OT}_S(e,f)\), so the product
again has the form \(p(u)k(v)\).  Similarly,
\[
 (p(u)k(v))^{-1}=k(v)^{-1}p(u)^{-1}
\]
can be put back in the same form by one application of
\eqref{eq:OT-refactor}.  Hence the graph is a subgroup.

Conversely, suppose that the graph is stable under multiplication and
apply this to the universal product \(k(v)p(u)\).  Its first two
coordinates are \(\widetilde u=(1+fv)u\) and
\(\widetilde v=v(1+du)\), and their defining equations reduce to
\[
 \widetilde u^2-c\widetilde u=-cfuv,\qquad
 \widetilde v^2-e\widetilde v=-eduv.
\]
Since
\[
 S[u,v]/(u^2-cu,v^2-ev)
\]
is free on \(1,u,v,uv\), both coefficients vanish.  Thus \(cf=ed=0\).
\end{proof}

\begin{remark}\label{rem:OT-scope}
Proposition~\ref{prop:OT-matched} is not a classification of all
rank-four group schemes.  It is an exact classification within the
following natural ansatz: the underlying scheme is the product of two
order-two Oort--Tate schemes, the multiplicative coordinate is the
product of their two characters, and the graph is required to lie in
\(H_S\).  The equations \(cf=ed=0\) are precisely the two conditions
that allow the factors to pass through one another.  The counterexample
first arises from this construction over \(D\), and then leaves the
product ansatz over \(R\) by modifying one of the quadratic equations.
\end{remark}

We now apply Proposition~\ref{prop:OT-matched} to \(S=D\) with
\begin{equation}\label{eq:our-OT-parameters}
 (c,d,e,f)=(ab,a,a^2,b).
\end{equation}
Indeed,
\[
 (ab)a=a^2b=-2,
 \qquad
 (a^2)b=-2,
 \qquad
 a^2a=0,
 \qquad
 (ab)b=0
 \quad\text{in }D.
\]
Consequently,
\begin{equation}\label{eq:A0}
 A_0
 =
 D[U,V]/(U^2-abU,\ V^2-a^2V)
 =
 D\{1,U,V,UV\},
\end{equation}
and
\begin{equation}\label{eq:matched-product}
 j_0:G_0:=\Spec A_0\hookrightarrow H_D,
 \qquad
 (U,V)\longmapsto
 \bigl(U,V,(1+aU)(1+bV)\bigr)
\end{equation}
is a finite locally free closed subgroup of rank four.

For later use, we record its law explicitly.  One has
\[
 D=\Z[a,b]/(a^3,b^2,a^2b+2),
 \qquad
 D_{\mathrm{add}}
 \simeq
 (\Z/4)\{1\}
 \oplus
 (\Z/2)\{a,a^2,b,ab\}.
\]
For a \(D\)-algebra \(C\), put
\[
 \ell(u,v)=(1+au)(1+bv).
\]
Then
\begin{equation}\label{eq:downstairs-law}
 (u_1,v_1)\star(u_2,v_2)
 =
 \bigl(
  u_1+\ell(u_1,v_1)u_2,\,
  v_1\ell(u_2,v_2)+v_2
 \bigr),
\end{equation}
with inverse
\begin{equation}\label{eq:downstairs-inverse}
 (u,v)^{-1}
 =
 \bigl((1+bv)u,\,v(1+au)\bigr).
\end{equation}
Moreover,
\begin{equation}\label{eq:downstairs-square}
 [2]^\#(U)=bUV,
 \qquad
 [2]^\#(V)=aUV,
 \qquad
 [4]=\unit.
\end{equation}
Here is a direct derivation of these formulas.  Put
\[
 L=1+au,\qquad M=1+bv,\qquad \ell=LM.
\]
The two Oort--Tate relations give \(L^2=M^2=1\), so
\(\ell^{-1}=\ell\).  Formula \eqref{eq:downstairs-law} is the
restriction of \eqref{eq:H-law}.  The ambient inverse is
\((-\ell u,-v\ell,\ell)\).  Since
\[
 -Lu=-u-au^2=u,\qquad -vM=-v-bv^2=v,
\]
where \(u^2=abu\), \(v^2=a^2v\), and \(a^2b=-2\), this is precisely
\eqref{eq:downstairs-inverse}.  Finally,
\[
 (1+\ell)u=buv,\qquad v(1+\ell)=auv.
\]
Indeed, these identities follow by expanding \(\ell=1+au+bv+abuv\)
and reducing the two quadratic relations; the terms containing \(2b\)
vanish in \(D\) because \(b^2=0\) and \(a^2b=-2\).
The power formula \eqref{eq:H-power} now gives the first two identities
of \eqref{eq:downstairs-square}.  Since \(\ell^2=1\) and
\(2(1+\ell)=0\) in \(A_0\), the same power formula gives
\([4]=\unit\).

The two Oort--Tate factors are not normal: ambient conjugation gives
\begin{align}
 p(u)k(v)p(u)^{-1}
   &=\bigl(buv,\,v(1+au)\bigr),
   \label{eq:conj-K}\\
 k(v)p(u)k(v)^{-1}
   &=\bigl((1+bv)u,\,auv\bigr).
   \label{eq:conj-P}
\end{align}
These identities follow by multiplying in \(H_D\) and reducing
\(u^2=abu\) and \(v^2=a^2v\).
Since \(aUV\) and \(bUV\) are nonzero in the free algebra
\eqref{eq:A0}, \(G_0\) is a genuine matched product rather than a
semidirect product with the displayed factors.  Appendix
\ref{app:no-semidirect} rules out any other rank-two semidirect
presentation.
\subsection{The corrected lift}

The same Oort--Tate parameters \eqref{eq:our-OT-parameters} no longer
satisfy the cross conditions over \(R\): one has
\[
 ed=a^3=0,
 \qquad
 cf=(ab)b=ab^2=aq\ne0.
\]
Thus the coefficientwise lift of the matched product \(G_0\) need not
be closed under multiplication.  The following identity measures the
failure and its correction.

\begin{lemma}[The closure identity]\label{lem:closure}
For \(t\in Q\), put
\[
 B_t
 =
 R[U,V]/(U^2-abU+tV,\ V^2-a^2V),
 \qquad
 \lambda=(1+aU)(1+bV).
\]
Then \(B_t\) is free over \(R\) on \(1,U,V,UV\).  In
\(B_t\ot_R B_t\), set
\[
 U_*=U_1+\lambda_1U_2,
 \qquad
 V_*=V_1\lambda_2+V_2.
\]
Then
\begin{equation}\label{eq:closure-defect}
 \begin{pmatrix}
  U_*^2-abU_*+tV_*\\
  V_*^2-a^2V_*\\
  (1+aU_*)(1+bV_*)-\lambda_1\lambda_2
 \end{pmatrix}
 =
 \begin{pmatrix}
  a(q-t)V_1U_2\\0\\0
 \end{pmatrix}.
\end{equation}
Moreover, \(\lambda\) is a unit in \(B_t\).
\end{lemma}

The proof is the elementary reduction recorded in
Appendix~\ref{app:closure}.  The three entries in
\eqref{eq:closure-defect} are simply the three defining equations
evaluated on the ambient product of two points.

For \(t=0\), the coefficientwise lift has the single nonzero defect
\[
 aqV_1U_2.
\]
It is nonzero because \(aq\ne0\) and
\(B_0\ot_RB_0\) is free on the sixteen tensor monomials.
For \(t=q\), all three entries in \eqref{eq:closure-defect} vanish.
Thus we replace the first Oort--Tate equation by
\begin{equation}\label{eq:corrected-equation}
 U^2-abU=0
 \qquad\rightsquigarrow\qquad
 U^2-abU=-qV.
\end{equation}
This gives precisely the algebra \(A\) of Theorem~\ref{thm:main}.

More explicitly, first quotient \(R[V]\) by the monic polynomial
\(V^2-a^2V\), and then adjoin \(U\) subject to the monic polynomial
\(U^2-abU+qV\).  It follows successively that
\begin{equation}\label{eq:A-basis}
 A=R\{1,U,V,UV\}
\end{equation}
is finite free of rank four.  Since \(\lambda\) is a unit, the map
\[
 j:G=\Spec A\longrightarrow H_R,
 \qquad
 (U,V)\longmapsto(U,V,\lambda)
\]
has a surjective comorphism: \(u,v,z,z^{-1}\) map to
\(U,V,\lambda,\lambda^{-1}\), respectively.  It is therefore a closed
immersion.  Equation \eqref{eq:closure-defect} with
\(t=q\) says that ambient multiplication restricts to \(G\), and the
ambient unit lies on \(G\).  Hence \(G\) is a finite locally free monoid
scheme; its associativity and unit laws are inherited from \(H_R\).

Modulo \(\mathfrak m=(a,b)\), one has
\[
 U^2=V^2=0,
 \qquad
 \lambda=1,
\]
and the two coordinates are primitive.  The closed fiber is therefore
\(\alpha_2\times\alpha_2\), in particular a group scheme.
Lemma~\ref{lem:monoid-to-group} now shows immediately that \(G\) is a
group scheme.  Dually, \(A\) is a Hopf algebra with
\[
 \eps(U)=\eps(V)=0,\qquad
 \Delta(U)=U\ot1+\lambda\ot U,
 \qquad
 \Delta(V)=V\ot\lambda+1\ot V,
 \qquad
 \Delta(\lambda)=\lambda\ot\lambda.
\]

This exhibits the counterexample as a small modification of the
Oort--Tate matched product.  The failed cross condition is the
coefficient \(aq\); adding \(qV\) to the first quadratic removes exactly
that defect.  The modification cannot be avoided by choosing a different
lift of \(b\): every lift \(b+x\), with \(x\in Q\), satisfies
\[
 (b+x)^2=b^2=q,
\]
because \(bQ=Q^2=0\).

\section{The fourth-power word}
\label{sec:power}

There are two conditions for the fourth power of a point of \(G\) to be
the unit.  Its multiplicative coordinate must satisfy
\(\lambda^4=1\), and its two translation coordinates must be killed by
\(N_4(\lambda)\).  The first condition does not imply the second over a
nonreduced ring: from
\[
 \lambda^4-1=(\lambda-1)N_4(\lambda)
\]
one cannot cancel \(\lambda-1\).  In fact, the whole counterexample is
concentrated in the small geometric sum \(N_4(\lambda)\).

Set
\[
 \theta=\lambda-1=aU+bV+abUV.
\]
A short reduction by the two quadrics gives
\begin{equation}\label{eq:theta-identities}
 \lambda^2=1+2bV,\qquad
 2\theta=2bV,\qquad
 \theta^2=0.
\end{equation}
For example, if \(L=1+aU\) and \(M=1+bV\), then
\[
 M^2=1,\qquad
 L^2=1+2bV,
\]
using \(U^2=abU-qV\), \(V^2=a^2V\), and
\eqref{eq:base-identities}.  Hence \(\lambda^2=L^2M^2=1+2bV\).
The middle identity in \eqref{eq:theta-identities} is immediate from
the displayed formula for \(\theta\) and \(2a=0\).  Substituting it in
\(\lambda^2=1+2\theta+\theta^2\) then gives \(\theta^2=0\).  In
particular, \(\lambda\) is a unit;
in fact \(\lambda^{-1}=1-\theta\).

Since \(4=0\) and \(\theta^2=0\),
\begin{equation}\label{eq:norms}
 N_4(\lambda)=2\theta=2bV,\qquad
 \lambda^4=1,\qquad
 N_8(\lambda)=N_4(\lambda)(1+\lambda^4)=0.
\end{equation}
Conceptually, because \(\lambda=1+\theta\) and \(\theta^2=0\), one has
\[
 N_4(\lambda)
 =1+(1+\theta)+(1+2\theta)+(1+3\theta)
 =4+6\theta=2\theta.
\]
Thus \(N_4(\lambda)\) need not vanish even though \(4=0\) and
\(\lambda^4=1\): its value here is \(2\theta=2bV\).  Multiplication by
the translation coordinate \(U\) gives the nonzero value \(2bUV\).
The ambient matrix formula \eqref{eq:H-power} now gives
\[
 [4]^\#(U)=2bVU=2bUV.
\]
This is nonzero by Lemma~\ref{lem:base} and the free basis
\eqref{eq:A-basis}.  On the other coordinate,
\[
 [4]^\#(V)=2bV^2=2a^2bV=-4V=0.
\]
At eight, both \(U\) and \(V\) map to zero and \(\lambda\) maps to one.
Therefore
\[
 [8]=\unit.
\]

Finally, \(G\) is noncommutative.  Indeed,
\[
 \Delta(U)-\tau\Delta(U)
 =(\lambda-1)\ot U-U\ot(\lambda-1),
\]
and the coefficient of \(V\ot U\) is the nonzero element \(b\in R\).
More precisely, this is its coefficient relative to the free tensor
basis coming from \eqref{eq:A-basis}, so it cannot be canceled by the
other displayed terms.

Lemma~\ref{lem:base} proves assertion (1) of
Theorem~\ref{thm:main}.  The successive monic presentation
\eqref{eq:A-basis} proves (2), and the closed-subgroup construction
together with Lemma~\ref{lem:monoid-to-group} proves (3).  The preceding
power calculation proves (4), including the nonvanishing of
\([4]^\#(U)\), while the closed-fiber and noncocommutativity calculations
prove (5).  This completes the proof of Theorem~\ref{thm:main}.

\appendix

\section{The hand calculation}\label{app:closure}

We prove Lemma~\ref{lem:closure}.  Fix \(t\in Q\).  The two successive
monic relations show first that \(B_t\) is free on \(1,U,V,UV\).
Work in \(B_t\ot_R B_t\).  For any \(x\in B_t\), write
\[
 x_1=x\ot1,\qquad x_2=1\ot x.
\]
Thus \(U_1,V_1\) belong to the first tensor factor and \(U_2,V_2\) to
the second.  In particular,
\[
 V_1U_2=(V\ot1)(1\ot U)
\]
is the mixed term that will measure the interaction between the two
points.  Scalars from \(R\), such as \(t\) and \(q\), are not
subscripted.

For \(i=1,2\), define
\begin{equation}\label{eq:tensor-characters}
 L_i=1+aU_i,\qquad
 M_i=1+bV_i,\qquad
 \lambda_i=L_iM_i.
\end{equation}
The coordinates of the ambient product are
\begin{equation}\label{eq:star-coordinates}
 U_*=U_1+\lambda_1U_2,
 \qquad
 V_*=V_1\lambda_2+V_2.
\end{equation}
We shall substitute these coordinates into the \(V\)-quadratic, the
\(U\)-quadratic, and the graph equation, in that order.

\smallskip
\noindent\emph{Preliminary reductions in each tensor factor.}
The base identities and the hypothesis \(t\in Q\) give
\[
 2a=4=0,\qquad
 t^2=bt=2t=0,\qquad
 a^2q=-2b.
\]
Inside either tensor factor,
\[
 U_i^2=abU_i-tV_i,
 \qquad
 V_i^2=a^2V_i.
\]
Consequently,
\[
 M_i^2
 =1+(2b+a^2q)V_i
 =1
\]
and
\[
 L_i^2
 =1+2aU_i+a^2U_i^2
 =1-a^2tV_i.
\]
Multiplying these two identities, and making three further immediate
reductions, gives
\begin{equation}\label{eq:auxiliary}
\begin{gathered}
 M_i^2=1,\qquad
 \lambda_i^2=1-a^2tV_i,\qquad
 a^2(\lambda_i-1)=-2V_i,\\
 t(\lambda_i-1)=atU_i,\qquad
 t(\lambda_i^2-1)=0,\\
 2\lambda_iU_i
   +ab(\lambda_i^2-\lambda_i)
   =-aqV_i
\end{gathered}
\qquad (i=1,2).
\end{equation}
For example, the middle identities follow directly from
\[
 \lambda_i-1=aU_i+bV_i+abU_iV_i
\]
and \(bt=t^2=0\).  For completeness, the two terms in the last
identity of \eqref{eq:auxiliary} reduce separately to
\[
 2\lambda_iU_i=2U_i+2bU_iV_i
\]
and
\[
 ab(\lambda_i^2-\lambda_i)
   =2U_i-aqV_i+2bU_iV_i.
\]
Their sum is \(-aqV_i\), since \(4=0\).

\smallskip
\noindent\emph{The \(V\)-quadratic.}
Expanding the second formula in \eqref{eq:star-coordinates} and using
\(V_i^2=a^2V_i\), we obtain
\begin{align*}
 V_*^2-a^2V_*
 &=V_1\lambda_2
   \bigl(a^2(\lambda_2-1)+2V_2\bigr)\\
 &=0.
\end{align*}
The last equality is exactly the third identity in
\eqref{eq:auxiliary}, applied in the second tensor factor.

\smallskip
\noindent\emph{The \(U\)-quadratic.}
Expanding the first formula in \eqref{eq:star-coordinates} and using
\(U_i^2=abU_i-tV_i\) in both factors gives
\begin{align*}
 U_*^2-abU_*+tV_*
 &=tV_1(\lambda_2-1)\\
 &\quad+
 \bigl(
  2\lambda_1U_1
  +ab(\lambda_1^2-\lambda_1)
 \bigr)U_2
 +t(1-\lambda_1^2)V_2.
\end{align*}
By \eqref{eq:auxiliary}, the three terms on the right are respectively
\[
 atV_1U_2,\qquad
 -aqV_1U_2,\qquad
 0.
\]
It follows that
\[
 U_*^2-abU_*+tV_*
 =a(t-q)V_1U_2
 =a(q-t)V_1U_2.
\]
The last equality uses \(2a=0\), so that \(-a=a\).

\smallskip
\noindent\emph{The graph equation.}
Set
\[
 \alpha=aU_2,\qquad
 \beta=bV_1,\qquad
 E=\alpha\beta.
\]
The subscripts are important: \(\alpha\) comes from the second tensor
factor and \(\beta\) from the first.  From
\eqref{eq:tensor-characters} and \eqref{eq:star-coordinates},
\[
 1+aU_*=L_1(1+M_1\alpha),
 \qquad
 1+bV_*=M_2(1+\beta L_2).
\]
Since \(M_1=1+\beta\) and \(L_2=1+\alpha\), while
\(\lambda_1\lambda_2=L_1M_2M_1L_2\), subtraction gives
\begin{align}
 &(1+aU_*)(1+bV_*)-\lambda_1\lambda_2\notag\\
 &\quad=
 L_1M_2
 \bigl(
  (1+M_1\alpha)(1+\beta L_2)-M_1L_2
 \bigr)\notag\\
 &\quad=
 L_1M_2E(2+\alpha+\beta+E).
 \label{eq:graph-factorization}
\end{align}
The last factorization is now an ordinary polynomial identity in
\(\alpha\) and \(\beta\).  It remains to see that every term in its
parenthesis is annihilated by \(E\).  First, \(2E=0\).  The quadratic
relations also give
\[
 \beta^2
 =b^2V_1^2
 =a^2qV_1
 =-2\beta,
 \qquad
 \alpha^2
 =a^2U_2^2
 =-a^2tV_2.
\]
Therefore
\[
 E\beta=-2E=0,
 \qquad
 E\alpha=-a^2btV_1V_2=0,
 \qquad
 E^2=0,
\]
where the middle equality uses \(bt=0\).  The right side of
\eqref{eq:graph-factorization} is zero, so the graph equation is
preserved.

We have now proved that the three defects are
\[
 \bigl(a(q-t)V_1U_2,\,0,\,0\bigr).
\]
Finally, in one copy of \(B_t\), the same preliminary calculation gives
\[
 \lambda^2=1-a^2tV.
\]
Since \((a^2tV)^2=0\), the character \(\lambda\) is a unit; explicitly,
\[
 \lambda^{-1}=\lambda(1+a^2tV).
\]
This proves every assertion of Lemma~\ref{lem:closure}.
\section{No semidirect presentation downstairs}
\label{app:no-semidirect}

The conjugation formulas \eqref{eq:conj-K}--\eqref{eq:conj-P} show
that the two displayed factors of \(G_0\) are not normal.  The following
strengthening rules out a different rank-two semidirect presentation.

\begin{corollary}\label{cor:no-semidirect}
The group scheme \(G_0\) has no finite locally free normal subgroup
scheme of rank two.  Consequently it is not isomorphic over \(D\) to a
nontrivial semidirect product of finite locally free group schemes.
\end{corollary}

\begin{proof}
Suppose that \(N\subset G_0\) is finite locally free of rank two.  On
the closed fiber, \(\mathcal O(N_{\F_2})\) is a two-dimensional quotient
of the local algebra
\[
 \F_2[U,V]/(U^2,V^2).
\]
As a quotient of this local algebra, it is local with residue field
\(\F_2\); its one-dimensional maximal ideal has square zero.  It is
therefore \(\F_2[T]/(T^2)\), and at least one of the images of \(U\) and
\(V\) is nonzero.  Hence at least one of the two scheme projections
\[
 N_{\F_2}\longrightarrow P_{\F_2},
 \qquad
 N_{\F_2}\longrightarrow K_{\F_2}
\]
is an isomorphism.  By Nakayama's lemma, the corresponding projection
over \(D\) is an isomorphism: on coordinate rings it is a map between
free \(D\)-modules of rank two whose reduction is invertible.

For the normality test we use the following identity in the ambient
group:
\begin{equation}\label{eq:ambient-conjugation}
 (x,y,z)(r,s,w)(x,y,z)^{-1}
 =
 \bigl(zr+(1-w)x,\ (s+y(w-1))z^{-1},\ w\bigr).
\end{equation}

Assume first that \(N\to P\) is an isomorphism.  In the product scheme
\(P\times_DK\), the subgroup \(N\) is then the graph
\(v=cu\) of a pointed scheme map, for some \(c\in D\): indeed, the
augmentation ideal of \(\mathcal O(P)=D\oplus DU\) is the principal
module \(DU\).  Along this graph,
\[
 \ell=1+(a+bc)u.
\]
Conjugation by the universal element \(k(y)\) gives
\begin{equation}\label{eq:normal-graph-P}
 V\bigl(k(y)n(u)k(y)^{-1}\bigr)
 -cU\bigl(k(y)n(u)k(y)^{-1}\bigr)
 =(a+bc)uy.
\end{equation}
If \(N\) were normal, the left side would vanish.  The algebra
\[
 D[u,y]/(u^2-abu,\ y^2-a^2y)
\]
is free on \(1,u,y,uy\), so \eqref{eq:normal-graph-P} would force
\(a+bc=0\).  This is impossible because \(a\notin bD\).

If instead \(N\to K\) is an isomorphism, write its graph as
\(u=cv\); here the same argument uses the principal augmentation ideal
\(DV\subset\mathcal O(K)\).  Conjugation by the universal \(p(x)\)
similarly gives
\begin{equation}\label{eq:normal-graph-K}
 U\bigl(p(x)n(v)p(x)^{-1}\bigr)
 -cV\bigl(p(x)n(v)p(x)^{-1}\bigr)
 =(b+ac)xv.
\end{equation}
Normality would force \(b+ac=0\), impossible because
\(b\notin aD\).  Indeed, the displayed additive decomposition of \(D\)
gives
\[
 bD\subset\langle2,b,ab\rangle,
 \qquad
 aD\subset\langle2,a,a^2,ab\rangle,
\]
which makes \(a\notin bD\) and \(b\notin aD\) immediate.  Thus no such
normal subgroup exists.

In a nontrivial semidirect decomposition of a rank-four group scheme
over the local ring \(D\), both factors have rank two and one is normal.
The first assertion therefore implies the second.
\end{proof}

\section{Companion computer audit}
\label{app:audit}

The proof above is independent of computer algebra.  For reproducibility,
the companion file
\begin{center}
\path{m2/verify_rank4_affine_article_20260712.m2}
\end{center}
performs exact strong Gr\"obner reductions over \(\Z\).  From the project
root, run
\begin{verbatim}
M2 --script m2/verify_rank4_affine_article_20260712.m2
\end{verbatim}
The script was tested with Macaulay2 1.25.11 and is organized around the
claims of the paper.  The article-specific blocks check:
\begin{itemize}
\item the base identities, the survival of \(q,aq,a^2q=-2b\), and the
normal form of \(2bUV\);
\item the exact three-component closure identity
\eqref{eq:closure-defect}, including the coefficientwise and corrected
lifts;
\item preservation of both quadrics, group-likeness of \(\lambda\),
coassociativity, the counit, and an explicit antipode satisfying both
convolution identities;
\item the explicit law, inverse, mutual actions, and conjugation formulas
for \(G_0\), including the nonnormality of \(P\) and \(K\), and the
graph identities used in Appendix~\ref{app:no-semidirect};
\item noncocommutativity and the special fiber
\(\alpha_2\times\alpha_2\);
\item the identities \(\theta^2=0\), \(2\theta=2bV\),
\[
 ([4]^\#U,[4]^\#V,[8]^\#U,[8]^\#V)=(2bUV,0,0,0);
\]
\item the general matched-product identities of
Proposition~\ref{prop:OT-matched} and the nonsplitting identity for the
chosen lift.
\end{itemize}
Freeness in \eqref{eq:A-basis}, the base complete-intersection assertion,
and the existence of the group law from the ambient group are proved
formally in the text; the script audits every polynomial identity used
in those arguments.  It also retains supplementary consistency checks
in standard semidirect-product coordinates and in a universal
five-parameter presentation.

\begin{thebibliography}{99}

\bibitem{SGA3I}
M.~Demazure and A.~Grothendieck (eds.),
\emph{Sch\'emas en groupes I: Propri\'et\'es g\'en\'erales des sch\'emas
en groupes (SGA 3)},
Lecture Notes in Mathematics, vol.~151, Springer-Verlag, 1970.

\bibitem{SGA3II}
M.~Demazure and A.~Grothendieck (eds.),
\emph{Sch\'emas en groupes II: Groupes de type multiplicatif, et structure
des sch\'emas en groupes g\'en\'eraux (SGA 3)},
Lecture Notes in Mathematics, vol.~152, Springer-Verlag, 1970.

\bibitem{SGA3VIIIannotated}
A.~Grothendieck,
\emph{Expos\'e VIII: Groupes diagonalisables},
revised and annotated online edition, ed. P.~Gille and P.~Polo,
version 1.1, 8 November 2009,
\href{https://wstein.org/sga/SGA3/Exp8-8nov09.pdf}
{available online}.

\bibitem{OortTate}
J.~Tate and F.~Oort,
\emph{Group schemes of prime order},
Ann. Sci. \'Ecole Norm. Sup. (4) \textbf{3} (1970), no.~1, 1--21,
\href{https://www.numdam.org/item/ASENS_1970_4_3_1_1_0/}
{Numdam}.

\bibitem{Schoof2001}
R.~Schoof,
\emph{Finite flat group schemes over local Artin rings},
Compos. Math. \textbf{128} (2001), no.~1, 1--15,
\href{https://doi.org/10.1023/A:1017560215203}
{doi:10.1023/A:1017560215203}.

\bibitem{SchoofSurvey}
R.~Schoof,
\emph{Is a finite locally free group scheme killed by its order?},
in \emph{Open Problems in Arithmetic Algebraic Geometry},
Advanced Lectures in Mathematics, vol.~46, International Press, 2019,
1--7,
\href{https://www.mat.uniroma2.it/~schoof/schoof_oortAAG.pdf}
{available online}.

\bibitem{Tate1997}
J.~Tate,
\emph{Finite flat group schemes},
in \emph{Modular Forms and Fermat's Last Theorem},
G.~Cornell, J.~H.~Silverman, and G.~Stevens (eds.),
Springer, New York, 1997, 121--154,
\href{https://doi.org/10.1007/978-1-4612-1974-3_5}
{doi:10.1007/978-1-4612-1974-3\_5}.

\bibitem{Torti}
E.~Torti,
\emph{Lagrange's theorem for a class of finite flat group schemes over
local Artin rings},
Doc. Math. (2025), online first,
\href{https://doi.org/10.4171/DM/1055}{doi:10.4171/DM/1055}.

\end{thebibliography}

\end{document}
